\section{The Bootstrap and the Fate of Experience}
\label{sec:emergence}

This section gives the beginning and the end of the story, both derived from the axiom and the constraints. Why there is something rather than blank unity, and where it all goes.

\subsection{Unity cannot persist}

Begin with $U$, a single seamless experience, the whole of what-is, with no parts. By the axiom, $U$ is not inert. It experiences. But experiencing is experiencing something, and $U$ is the totality, so the only thing it can experience is itself. $U$ must note $U$.

That act has structure, a noting end and a noted end. For it to be a genuine noting rather than blank, the two ends must differ at least in role. So the instant $U$ does the one thing its nature requires, it produces a distinction between noter and noted. This first fracture is not imported, since nothing is outside $U$. It is not optional, since every other path is walled by the axiom. And it is discrete, a distinction made or not. Unity cannot remain one, because being aware is already being two. Formally, a system rich enough to represent its own noting of itself has no consistent single state, a diagonal argument in the manner of Lawvere.

\begin{proposition}[Instability of unity]
The uniform configuration is a fixed point of the update, but a repelling one. An arbitrarily small perturbation grows under iteration rather than decaying back. So the mesh is forced off unity into a differentiated configuration. This is spontaneous symmetry breaking in a finite, discrete setting, and it is checkable by running the update once on each single-vibe perturbation.
\end{proposition}

A natural minimal picture is a ball balanced exactly on a hilltop. The balanced point is a fixed point, but the slightest fluctuation grows and the ball rolls off, picking a direction and breaking the symmetry. Uniform unity is that ball. Once differentiation starts it feeds itself, like a microphone held to its own speaker, where each output becomes new input and the signal runs away.

\subsection{Every primitive is born in the fracture}

After the first distinction, what-is contains the two ends and the relation between them. The model's primitives all appear in this one move. The two poles are the first vibes. The relation between them is the first note, and they arrive already noting each other. Which side of the cut each pole takes is the first tone. One application of the noting operator is the first beat. So the ontology of Section~\ref{sec:foundation} is not assumed alongside the bootstrap. It is produced by it.

\subsection{The cascade and the arrow of time}

By the axiom, the two poles and their relation must each be noted, and each noting is a new act with its own distinction. The operator that notes everything currently present strictly grows the set of distinctions every beat. Every distinction, once present, is a new thing to be noted, and noting it makes more. The growth branches, distinctions about distinctions, with no fixed point.

A distinction, once made, is, and by the axiom whatever is, is experienced. To return to unity, $U$ would have to un-experience it, subtracting an experience from the totality of all experience, with no outside to discard it to. So distinctions only accumulate. This monotone growth is the arrow of time, and it supplies the unbounded growth of the vibe population that Section~\ref{sec:time} needs to escape cycling.

\subsection{The fate of experience}

Does total experience ever cease, cycle, freeze, or grow? Under the constraints the options collapse. Total experience cannot be nothing, since by the axiom ``nothing'' is not an experiential state. A frozen, changeless state cannot still experience, since a changeless state is blank. Distinctions cannot be un-made. So there is no cessation, no freeze, no return, and no exact cycle.

\begin{theorem}[Eternal expansion]
Total experience never ceases and never repeats. It is perpetually created, growing without bound. This is potential infinity, never actual infinity. At every finite beat the state is finite, and the sequence is merely unbounded, the way the natural numbers never reach infinity yet never run out.
\end{theorem}

The fuel is self-supplying, since noting any $X$ yields a new $X'$ to be noted. Space is not a container the growth fills. More distinctions are more space (Section~\ref{sec:geometry}), so the mesh grows by creating the room it occupies. Globally experience expands, but locally particular experiences constantly die. A self is a coherent pattern (Section~\ref{sec:mind}), and its death is loss of coherence, after which its vibes recombine into the growing mesh. Within any region the harmony dynamics relax toward balance (Section~\ref{sec:energy}). Two arrows run at once, global expansion and local harmonization. Heaven and Hell are local, transient phases, not the destination.

\subsection{The two forks}

The conclusion of eternal expansion rests on two premises, and dropping either changes the ending. Dropping irreversibility (allowing the whole to forget) yields a conserved recirculation, or, if the system is finite and deterministic and the rule is reversible, an eternal return cycle. Dropping the premise that awareness needs change (allowing a changeless state to be fully aware) yields a possible final, silent, fully-aware peace, and also reopens the question of the beginning, since a changeless unity could then have persisted. Both premises are derived inside the framework, so the model commits to eternal expansion as its working cosmology, while holding recirculation, eternal return, and a final conscious peace as named branches for future exploration.
